\sscp{PMP *q and *h in Wallacea}{13-01-04}
Even though there is a widespread impression of the loss of PMP
*q and *h in the Wallacean region,\footnote{
		As evidenced in various
		emails from other linguists sent to us.}
nevertheless it would
be a disservice to the data to make such a sweeping generalisation,
or to remove those proto-phonemes from reconstructions covering the whole region.
Some individual Wallacean subgroups,
such as \tsc{Sula-Buru} do show regular PMP *q/*h > {\0}.

However, in \srf{11-03-01} we showed that the \tsc{Seram-Tanimbar-Bomberai}
language Watubela retains PMP *q in all positions.
In \trf{04-15} (\prf{tab:04-15}), we showed that several \tsc{Timor-Babar} languages
in eastern Timor retain reflexes of medial *q.
\citet[82]{Edwards2021} also notes “reflexes of *q in some languages
in some forms [{\ldots}] indicate that a reflex of it probably
needs to be reconstructed in some forms for Proto-Timor-Babar.”
There are examples from Rote,
for example, of sporadic raising PMP *-aq > Proto-Rote **ə instead
of the expected **a \citep[151]{Edwards2021}.
In \tsc{Ambon-Seram} languages, retention of initial *q has been
proposed \citep[121ff]{Collins1983} as explaining some unexpected
consonant onsets on some nouns (discussed in \srf{13-02-02},
where alternate explanations are also given).
Also in \tsc{Ambon-Seram} languages,
retention of medial *q has been suggested as explaining why the
reflexes of historical secondary vowel sequences *aqi
and *aqu are often different from *ai,
*au, *a(C)i, and *a(C)u that result from the loss of other historical
medial consonants (discussed in \srf{10-02}, example 24).
And in the \tsc{Ambon-Seram} language Alune,
final *q{\#} > \it{kw-e{\#}} in several words (\srf{10-02}, example \ref{ex:10-22}).
So while PMP *q does not seem to have subgrouping value as a
whole: it is retained fully in Watubela,
its traces in other languages may explain conditioned reflexes
at some lower level nodes,
and there also seem to be lexically specific retentions scattered throughout the region.

\citet{Blust1981} argued for the retention of \tsc{PAn} *S (PMP
*h) in Soboyo, a \tsc{Taliabu} language (linked with \tsc{Celebic}; see \crf{ch:Tal}).
The normal reflex in Soboyo is PMP *h > {\0},
so its possible retention in word-initial
and word-medial positions reflects either an unconditioned split,
or lexically specific retentions.
The words which seem to reflect initial PMP *h = \it{h} in Soboyo
are perhaps better explained by the widespread phenomenon of
irregular C-onsets on vowel-initial roots (discussed in \srf{13-02-02}).
The one medial form that is deemed “irreproachable” \citep[25]{Blust1981},
is Soboyo \it{kuhi} ‘knock down with an upward blow,
as fruit from a tree,
etc.’ assigned to \tsc{PAn} *kuSit.
However, \tsc{PAn} *kuSit is not found in \citet{BlustTrussel2020},
but there is PWMP *kuhit ‘tap,
touch, or reach in with the finger’.
It is unclear if the Soboyo semantics of ‘knock down’ are actually
related to the glosses hovering around ‘reach in with the finger’.
And Soboyo retains final *t > \it{c} in some etyma (e.g.
*ma-bəRat > \it{m-fahac} ‘heavy’,
*əpat > \it{ŋ-haac} ‘four’,
PCM **səgət \it{> segec} ‘high tide’),
so its cognacy with the PWMP form is not “irreproachable” on
the basis of either form or meaning.
It is nevertheless one possible retention of PMP *h.

As with PMP *q,
\citet[41ff]{Collins1983} argues that retention of medial PMP
*h in some \tsc{Ambon-Seram} languages explains why the reflexes
of historical secondary vowel sequences *ahi
and *ahu are sometimes different from*ai,
*au, *a(C)i, and *a(C)u that result from the loss of other historical
medial consonants (discussed in \srf{10-02},
where we note that the historical coda may be the conditioning
factor in the vowel shift,
rather than a historical medial consonant).
Claims about the retention of PMP *h in the region are sporadic
and irregular, often lexically specific, and usually have other
reasonable explanations as alternatives.

Regardless, the loss of PMP *q
and *h is also found widely west
and north of the Blust line,
in many individual languages
and in some whole subgroups.
For example, *h > {\0} is one of the defining features of \tsc{Celebic}
languages \citep{Mead2003a}.
But because possible traces influence conditioned reflexes throughout
the region, more work is needed,
and care must be taken to not lightly dismiss possible retentions.
The possibility of PMP *h > {\0} as an innovation that unifies
Wallacean subgroups needs further study.
Even if this change can be accepted for all Wallacean subgroups,
its cross linguistic frequency,
and the fact that it is not exclusive means that its value for
subgrouping would be almost negligible.

